% ============================================================
% ch16.tex —— 第 16 章 求导工具箱：自己发明运算法则
% ============================================================
\chapter{求导工具箱：自己发明运算法则}

上一章的机器很好，但每换一个函数就要重新展开一次 $(t+h)^n$——手工作坊，效率太低。这一章把作坊升级成流水线：\textbf{四条法则}，覆盖本书所需的全部求导。每一条都现场发明，没有一条靠背——发明完顺手处死两个初中公式（顶点公式、切线方程），让你尝到流水线的甜头。

\section{法则〇：常数不动}

最简单的函数：$f(t) = c$（常数，图像一条水平线）。差商 $\dfrac{c - c}{h} = 0$，落脚点 $0$：
\[
\boxed{\;(c)' = 0\;}
\]
不变的东西变化率为零——图像水平，斜率为 $0$，几何与代数互相印证。这条"法则〇"看似废话，却是流水线的保险丝（后面所有法则都要用它打底）。

\section{法则一：幂法则——从数据到证明}

第 15 章已攒四票数据：$t^1\to1$，$t^2\to2t$，$t^3\to3t^2$，$t^4\to4t^3$，外加分数指数的两票 $\sqrt t\to\tfrac1{2\sqrt t}$（即 $(t^{1/2})' = \tfrac12t^{-1/2}$）、$\tfrac1t\to-\tfrac1{t^2}$（即 $(t^{-1})' = -t^{-2}$）。六票全部支持同一个猜想：
\[
\boxed{\;(t^n)' = n\,t^{n-1}\;}
\]
（指数落下来当系数，指数减一。对正整数、零、负数、分数指数都成立。）对正整数 $n$ 的证明只需一次通用展开——二项式展开的首两项：
\[
(t+h)^n = t^n + n\,t^{n-1}h + \binom{n}{2}t^{n-2}h^2 + \cdots + h^n .
\]
（$\binom n2 = \tfrac{n(n-1)}2$ 等是组合系数；$n = 2, 3, 4$ 时这个"首两项"模式你已经亲眼乘出来过。二项式系数的完整故事——杨辉三角——在第 22 章还有一次出场。）于是差商：
\[
\frac{(t+h)^n - t^n}{h} = \frac{n t^{n-1}h + \binom n2 t^{n-2}h^2 + \cdots + h^n}{h} = n t^{n-1} + \binom n2 t^{n-2}h + \cdots + h^{n-1} .
\]
除完之后，\textbf{除首项外每一项都挂着 $h$}——比赛里全体压成零，落脚点 $n t^{n-1}$。证毕。负指数与分数指数的证明需要指数函数的工具（第 19 章），先记账。

\section{法则二：和法则——免费的午餐}

\[
\boxed{\;(f + g)' = f' + g'\;}
\]
\begin{proof}
差商拆开即可：
\[
\frac{[f+g](t+h) - [f+g](t)}{h} = \frac{f(t+h)-f(t)}{h} + \frac{g(t+h)-g(t)}{h} \;\longrightarrow\; f' + g' .
\]
（两个分别收敛的量相加，和收敛到极限之和——第 14 章规则的日常应用。）
\end{proof}
\textbf{变化互不串门}：各自的增量独立结算。配合法则〇立刻得到使用率最高的一条：\textbf{多项式逐项求导}：
\[
\left(3t^5 - 2t + 100\right)' = 15t^4 - 2 + 0 = 15t^4 - 2 .
\]
（常数 $100$ 求导蒸发——这就是"平移不改变斜率"的代数版。）

\section{法则三：乘法法则——画一个长方形}

$(f\times g)'$ 等于什么？第一直觉是"$f'\times g'$"——\textbf{用数据击毙它}：取 $f = g = t$，左边 $(t\cdot t)' = (t^2)' = 2t$，右边 $1\times1 = 1$，不等。直觉错在哪？错在忽略了"两边同时在动"的交叉效应。正确的直觉来自几何：

设 $f, g$ 是一块长方形的两条边，面积 $fg$。时间从 $t$ 走到 $t+h$，两边各伸长 $\Delta f = f(t+h)-f(t)$、$\Delta g$。新增面积分三块：

\begin{center}
\begin{tikzpicture}[scale=1.0]
  \draw[fill=shenlan!12] (0,0) rectangle (3,2);
  \draw[fill=molv!25] (3,0) rectangle (3.5,2);
  \draw[fill=cheng!30] (0,2) rectangle (3,2.5);
  \draw[fill=shaohong!40] (3,2) rectangle (3.5,2.5);
  \draw (0,0) rectangle (3.5,2.5);
  \node at (1.5,1) {$f$};
  \node at (3.25,1) {$\Delta f$};
  \node at (1.5,2.25) {\scriptsize $g\,\Delta f$};
  \node at (3.25,2.25) {\scriptsize $\Delta f\,\Delta g$};
  \node[right] at (3.6,1) {\scriptsize $f\,\Delta g$};
  \draw[<->] (0,-0.35) -- (3,-0.35) node[midway,below] {\scriptsize $f$};
  \draw[<->] (3.02,-0.35) -- (3.5,-0.35) node[below] {\scriptsize};
  \draw[<->] (-0.35,0) -- (-0.35,2) node[midway,left] {\scriptsize $g$};
\end{tikzpicture}
\end{center}

右侧竖条面积 $g\,\Delta f$，顶部横条 $f\,\Delta g$，角落小块 $\Delta f\,\Delta g$。除以 $h$：
\[
\frac{\Delta(fg)}{h} = g\,\frac{\Delta f}{h} + f\,\frac{\Delta g}{h} + \frac{\Delta f\,\Delta g}{h} .
\]
前两项分别奔 $g f' + f g'$。第三项呢？$\Delta f$ 与 $h$ 同级小（$\Delta f \approx f'h$），$\Delta g$ 亦然，于是
\[
\frac{\Delta f\,\Delta g}{h} \approx \frac{f'g'\,h^2}{h} = f'g'\,h \;\longrightarrow\; 0
\]
——\textbf{角落小块是"平方级的小"，比 $h$ 小得多，比赛里蒸发}。（这与第 15 章"切线定义"里被挤掉的弯曲量是同一种小量。）所以：
\[
\boxed{\;(fg)' = f'g + fg'\;}
\]
检验两发。第一发：$(t\cdot t)' = 1\cdot t + t\cdot 1 = 2t$ \checkmark。第二发：$(t^2\cdot t^3)'$。乘法法则：$(t^2)'t^3 + t^2(t^3)' = 2t\cdot t^3 + t^2\cdot 3t^2 = 2t^4 + 3t^4 = 5t^4$；而把乘积先展开成 $t^5$ 再用幂法则，$(t^5)' = 5t^4$——两种路线同一答案 \checkmark。这道对账题的价值：多项式相乘容易错漏，法则保证结构正确。

\section{法则四：链式法则——穿盒术}

函数套函数怎么求导？比如 $y = (t^2+1)^{10}$：内层 $u = t^2+1$，外层 $y = u^{10}$。硬展开十次方是自杀，正确姿势是\textbf{穿盒}：

比喻：人民币换美元、美元换欧元，求"人民币兑欧元"的汇率——\textbf{两段汇率相乘}。变化率完全同理：$y$ 随 $u$ 变化的速度，乘以 $u$ 随 $t$ 变化的速度：
\[
\boxed{\;\frac{dy}{dt} = \frac{dy}{du}\cdot\frac{du}{dt}\;}
\]
莱布尼茨记号在这里显出威力：分式外观上"$du$"像被约掉——不是真的约去（它们不是分数，是记号），但\textbf{记忆结构与分数一致}，这正是莱布尼茨设计记号的苦心。

实战：$y = (t^2+1)^{10}$。外层对 $u$ 求导：$10u^9$；内层：$2t$。相乘：
\[
y' = 10(t^2+1)^9\cdot 2t = 20t\,(t^2+1)^9 .
\]
一行搞定十次方。再来两个：
\[
y = \sqrt{t^2+1}:\quad y' = \frac{1}{2\sqrt{t^2+1}}\cdot 2t = \frac{t}{\sqrt{t^2+1}} .
\]
\[
y = (3t+7)^5:\quad y' = 5(3t+7)^4\cdot 3 = 15(3t+7)^4 .
\]
（第二个请注意"\textbf{内层导数要乘上去}"：内层 $3t+7$ 的导数是 $3$，漏乘是初学者第一大错误——检验方法是取 $t$ 使内外层都好算，代回原函数的差商抽查一遍。）

\begin{cidian}
四法则的正式户口：\textbf{幂法则}（power rule）、\textbf{和法则}（sum rule）、\textbf{乘法法则}（product rule，莱布尼茨法则）、\textbf{链式法则}（chain rule）。加上法则〇（常数的导数为零），五件套在手，本书范围内\textbf{任何}函数的求导都是代数作业。
\end{cidian}

\section{当场回本：处死两个初中公式}

\paragraph{处死一：二次函数顶点公式。}初中人人背过：$y = ax^2 + bx + c$ 的顶点横坐标是 $x = -\dfrac{b}{2a}$。背的时候你信过它吗？现在它是三行流水线作业：

抛物线的顶点是"转弯点"——左侧上升、右侧下降（或反之），顶点处升降恰好停顿，即\textbf{变化率为零}：
\[
y' = 2ax + b = 0 \quad\Longrightarrow\quad x = -\frac{b}{2a} .
\]
\textbf{公式是推出来的，不是背出来的}。而且你白得一个大礼包：这个方法找\textbf{任何}函数的最大最小值（第 18 章的物理题直接受益）——"导数为零处"是所有优化问题的第一嫌疑人。初中老师没骗你，但只给了鱼，没给渔。

\paragraph{处死二：切线方程。}求过 $y = x^2$ 上点 $(1,1)$ 的切线。切线是直线，设 $y = kx + m$。斜率 $k$ 就是该点的导数：$k = y'(1) = 2\times1 = 2$。代入点 $(1,1)$：$1 = 2 + m$，$m = -1$。切线：
\[
y = 2x - 1 .
\]
三行完事。初中这叫"超纲"，现在它只是流水线上的标准件。（顺带检验：这条切线与抛物线的公共点——联立 $x^2 = 2x - 1$，即 $(x-1)^2 = 0$，重根 $x=1$，恰好只有一个"贴点"——第 29、30 章的"相切 $=$ 重根"在此预演。）

\begin{dongshou}
\begin{itemize}
  \yisi 求导（流水线作业）：$t^7$；$3t^5 - 2t + 100$；$6t^2 + \tfrac5t$（提示 $\tfrac5t = 5t^{-1}$）；$(2t+1)^6$（链式）。
  \yier 用乘法法则求 $(t^2+1)(t^3-2)$ 的导数，再把括号乘开逐项求导，两种方法对账。（展开式 $t^5 + t^3 - 2t^2 - 2$，导数 $5t^4 + 3t^2 - 4t$；乘法法则应给出同一答案—— $(2t)(t^3-2) + (t^2+1)(3t^2)$，展开核对。）
  \yier 证明 $(cf)' = c\,f'$（常数可提到导数外）：用乘法法则，$(cf)' = 0\cdot f + c f'$。（"常数倍数不动，只对函数求导"——流水线最常用的短路径。）
  \yier 用"顶点公式三行版"重推：$y = -2x^2 + 8x - 3$ 的顶点（$x = -\tfrac{8}{2\times(-2)} = 2$，$y = 5$），并用配方法 $-2(x-2)^2+5$ 对账。两种方法哪种是"渔"？
\end{itemize}
\end{dongshou}

\begin{shaonao}
\textbf{用导数给函数画像。}三次函数 $y = x^3 - 3x$。求导 $y' = 3x^2 - 3 = 3(x-1)(x+1)$，零点 $x = \pm1$。符号分析：$x < -1$ 时 $y' > 0$（上升）；$-1 < x < 1$ 时 $y' < 0$（下降）；$x > 1$ 时 $y' > 0$（上升）。于是函数的"地形图"全在手：$(-1, 2)$ 是山顶，$(1, -2)$ 是山谷。

回答三问：(1) \textbf{三次函数最多几个极值点？}——导数是二次的，至多两个实根，故至多两处"变向"。(2) 四次函数呢？（导数三次，至多三个极值点——\textbf{$n$ 次函数至多 $n-1$ 个极值点}，与"至多 $n$ 个根"（第 11 章）一线相承。）(3) 这个方法什么时候失灵？——导数无实根时（如 $y = x^3$，$y' = 3x^2$ 只在 $0$ 处为零但不变号，$0$ 不是极值点，是"拐点"）：\textbf{导数为零是极值的嫌疑人，不是定罪}，定罪要加"变号"证据。能想到第 (3) 问，你已经站在大学微积分课的门口。
\end{shaonao}

\begin{chengshi}
乘法法则证明中"$\Delta f$ 与 $h$ 同级小"的严格版是"$\Delta f = f'h + o(h)$"（高阶无穷小记号），链条完整后它就是"微分"的定义；链式法则的完整证明（含 $\Delta u = 0$ 的退化情形）同样放行大学。幂法则对一切实数指数成立的证明需要指数函数理论（第 19 章后可补）。这些地毯下都是平整的地板，不影响行走。
\end{chengshi}

工具箱满载。但导数只回答了"变化多快"；它的孪生兄弟还压在案上：\textbf{变化积攒起来一共多少？}下一章，发明积分。
